% TODO proof-reading
\subsection{シンプルな関数呼び出し}

\TT{Math.random()}は[0.0 \dots 1.0)の範囲で疑似乱数を返しますが、
何らかの理由で[0.0 \dots 0.5)の範囲で数値を返す関数を考案する必要があるとしましょう：

\begin{lstlisting}[style=customjava]
public class HalfRandom
{ 
	public static double f()
	{
		return Math.random()/2;
	}
}
\end{lstlisting}

\begin{lstlisting}[caption=定数プール]
...
   #2 = Methodref          #18.#19    //  java/lang/Math.random:()D
   #3 = Double             2.0d
...
  #12 = Utf8               ()D
...
  #18 = Class              #22        //  java/lang/Math
  #19 = NameAndType        #23:#12    //  random:()D
  #22 = Utf8               java/lang/Math
  #23 = Utf8               random
\end{lstlisting}

\begin{lstlisting}
  public static double f();
    flags: ACC_PUBLIC, ACC_STATIC
    Code:
      stack=4, locals=0, args_size=0
         0: invokestatic  #2          // Method java/lang/Math.random:()D
         3: ldc2_w        #3          // double 2.0d
         6: ddiv          
         7: dreturn       
\end{lstlisting}

\TT{invokestatic}は\TT{Math.random()}関数を呼び出し、結果を\ac{TOS}に残します。

その後、結果は2.0で割られて返されます。

しかし、関数名はどのようにエンコードされているのでしょうか？

それは\TT{Methodref}式を使用して定数プールにエンコードされています。

これはクラス名とメソッド名を定義します。

\TT{Methodref}の最初のフィールドは\TT{Class}式を指しており、これは順番に
通常のテキスト文字列（「java/lang/Math」）を指しています。

2番目の\TT{Methodref}式は\TT{NameAndType}式を指しており、これも
文字列への2つのリンクを持っています。

最初の文字列は「random」で、これはメソッドの名前です。

2番目の文字列は「()D」で、これは関数の型をエンコードしています。
これは\IT{double}値を返すことを意味します（したがって文字列の\IT{D}）。

これが次の方法です：
1) JVMが型の正しさのためにデータをチェックできる；
2) Java逆コンパイラがコンパイルされたクラスファイルからデータ型を復元できる。

今度は「Hello, world!」の例を試してみましょう：

\begin{lstlisting}[style=customjava]
public class HelloWorld
{
	public static void main(String[] args)
	{
		System.out.println("Hello, World");
	}
}
\end{lstlisting}

\begin{lstlisting}[caption=定数プール]
...
   #2 = Fieldref           #16.#17        //  java/lang/System.out:Ljava/io/PrintStream;
   #3 = String             #18            //  Hello, World
   #4 = Methodref          #19.#20        //  java/io/PrintStream.println:(Ljava/lang/String;)V
...
  #16 = Class              #23            //  java/lang/System
  #17 = NameAndType        #24:#25        //  out:Ljava/io/PrintStream;
  #18 = Utf8               Hello, World
  #19 = Class              #26            //  java/io/PrintStream
  #20 = NameAndType        #27:#28        //  println:(Ljava/lang/String;)V
...
  #23 = Utf8               java/lang/System
  #24 = Utf8               out
  #25 = Utf8               Ljava/io/PrintStream;
  #26 = Utf8               java/io/PrintStream
  #27 = Utf8               println
  #28 = Utf8               (Ljava/lang/String;)V
...
\end{lstlisting}

\begin{lstlisting}
  public static void main(java.lang.String[]);
    flags: ACC_PUBLIC, ACC_STATIC
    Code:
      stack=2, locals=1, args_size=1
         0: getstatic     #2        // Field java/lang/System.out:Ljava/io/PrintStream;
         3: ldc           #3        // String Hello, World
         5: invokevirtual #4        // Method java/io/PrintStream.println:(Ljava/lang/String;)V
         8: return        
\end{lstlisting}

オフセット3の\TT{ldc}は定数プール内の「Hello, World」文字列へのポインタを取り、
スタックにプッシュします。

Javaの世界では\IT{reference}と呼ばれますが、これはむしろポインタ、またはアドレスです

\footnote{C++でのポインタと\IT{reference}の違いについては：\myref{cpp_references}。}。

馴染みのある\TT{invokevirtual}命令は定数プールから\TT{println}
関数（またはメソッド）に関する情報を取得し、それを呼び出します。

ご存知のように、各データ型に対して複数の\TT{println}メソッドがあります。

私たちの場合は、\IT{String}データ型用の\TT{println}のバージョンです。

しかし、最初の\TT{getstatic}命令はどうでしょうか？

この命令はオブジェクト\TT{System.out}のフィールドの\IT{reference}（またはアドレス）を取り、
スタックにプッシュします。

この値は\TT{println}メソッドの\IT{this}ポインタのように機能します。

したがって、内部的には、\TT{println}メソッドは入力として2つの引数を取ります：
1) \IT{this}、つまりオブジェクトへのポインタ；
2) 「Hello, World」文字列のアドレス。

実際、\TT{println()}は初期化された\TT{System.out}オブジェクト内のメソッドとして呼び出されます。

便宜上、\TT{javap}ユーティリティはこのすべての情報をコメントに書き込みます。